% Questo file definisce lo stile che verrà applicato
% ad ogni pagina di contenuto
\documentclass[a4paper,11pt]{article}

\usepackage{ifthen}
\usepackage[
 a4paper,
 top=2.5cm,
 bottom=2.5cm,
 left=1.5cm,
 right=1.5cm,
 head=30pt
]{geometry}
\usepackage[italian]{babel}
\usepackage[utf8x]{inputenc}
\usepackage[T1]{fontenc}
\usepackage{fancyhdr}
\usepackage[colorlinks=true, urlcolor=black, citecolor=black, linkcolor=black]{hyperref}
\usepackage{tabularx}
\usepackage{multirow}
\usepackage{booktabs}
\usepackage{color}
\usepackage[dvipsnames]{xcolor}
\usepackage{graphicx}
\usepackage{eurosym}
\usepackage{amsmath}
\usepackage{relsize}
\usepackage{placeins}
\usepackage{ltablex}
\usepackage{float}

\usepackage[multidot]{grffile}
\usepackage{xcolor,colortbl}
\definecolor{lightblue}{HTML}{56B4E6}
\definecolor{blue}{HTML}{2953A1}
\definecolor{darkblue}{HTML}{1E396E}
\usepackage{longtable}

\usepackage[toc,page]{appendix}
\renewcommand\appendixtocname{Appendice}
\renewcommand{\appendixpagename}{Appendice}

\newcommand\pagenumberingnoreset[1]{\gdef\thepage{\csname @#1\endcsname\c@page}}

% Cambia il font 
\renewcommand*\rmdefault{qhv}

% ***STILE PAGINA***
\pagestyle{fancy}
\fancyhf{}
\setlength{\headheight}{1cm} 
% No indentazione paragrafo
\setlength{\parindent}{0pt}

% ***INTESTAZIONE***
\newcommand\textline[4][t]{%
  \noindent\parbox[#1]{.333\textwidth}{\raisebox{-0.40\height}{#2}}%
  \parbox[#1]{.333\textwidth}{\centering #3}%
  \parbox[#1]{.333\textwidth}{\raggedleft #4}%
}

\lhead{
	\textline[t]{\includegraphics[width=1cm, keepaspectratio=true]{../../../Template/Logo/Logo.png}}{\progettoShort}{\documento}
}

\renewcommand{\headrulewidth}{0.4pt}  %Linea sotto l'intestazione

% ***PIÈ DI PAGINA***
\lfoot{\textit{\gruppoLink}\\ \footnotesize{\email}}
\rfoot{\thepage} %per le prime pagine: mostra solo il numero romano
\cfoot{}
\renewcommand{\footrulewidth}{0.4pt}   %Linea sopra il piè di pagina


% Ridefinisce command \paragraph{} andando a capo ogni dopo la parola dentro le parentesi ed ha la possibiltà di enumerazione fino a n cifre modificando il numero dentro "secnumdepth"
\usepackage{titlesec}

\setcounter{secnumdepth}{7}
\setcounter{tocdepth}{7}


% Visualizza paragraph come una section
\titleformat{\paragraph}{\normalfont\normalsize\bfseries}{\theparagraph}{1em}{}
\titlespacing*{\paragraph}{0pt}{3.25ex plus 1ex minus .2ex}{1.5ex plus .2ex}

\titleformat{\subparagraph}{\normalfont\normalsize\bfseries}{\thesubparagraph}{1em}{}
\titlespacing*{\subparagraph}{0pt}{3.25ex plus 1ex minus .2ex}{1.5ex plus .2ex}

\makeatletter
\newcounter{subsubparagraph}[subparagraph]
\renewcommand\thesubsubparagraph{%
  \thesubparagraph.\@arabic\c@subsubparagraph}
\newcommand\subsubparagraph{%
  \@startsection{subsubparagraph}    % counter
    {6}                              % level
    {\parindent}                     % indent
    {3.25ex \@plus 1ex \@minus .2ex} % beforeskip
    {0.75em}                           % afterskip
    {\normalfont\normalsize\bfseries}}
\newcommand\l@subsubparagraph{\@dottedtocline{6}{13em}{5.5em}} %gestione dell'indice
\newcommand{\subsubparagraphmark}[1]{}
\makeatother

\makeatletter
\newcounter{subsubsubparagraph}[subsubparagraph]
\renewcommand\thesubsubsubparagraph{%
  \thesubsubparagraph.\@arabic\c@subsubsubparagraph}
\newcommand\subsubsubparagraph{%
  \@startsection{subsubsubparagraph}    % counter
    {7}                              % level
    {\parindent}                     % indent
    {3.25ex \@plus 1ex \@minus .2ex} % beforeskip
    {0.75em}                           % afterskip
    {\normalfont\normalsize\bfseries}}
\newcommand\l@subsubsubparagraph{\@dottedtocline{7}{16em}{6.5em}} %gestione dell'indice
\newcommand{\subsubsubparagraphmark}[1]{}
\makeatother

%Generali
\newcommand{\capitolato}{C5 - Monolith: An interactive bubble provider}
\newcommand{\progettoShort}{Monolith}
\newcommand{\progetto}{Monolith: An interactive bubble provider}
\newcommand{\gruppo}{NPE Developers}
\newcommand{\gruppoLink}{\href{https://gitlab.com/npe-developers}{NpeDevelopers}}
\newcommand{\email}{\href{mailto:npe.developers@gmail.com}{\textcolor{blue}{npe.developers@gmail.com}}}
\newcommand{\password}{NP3Devel0pers}
\newcommand{\myincludegraphics}[2][]{%
	\setbox0=\hbox{\phantom{X}}%
	\vtop{
		\hbox{\phantom{X}}
		\vskip-\ht0
		\hbox{\includegraphics[#1]{#2}}}
}




%Componenti del gruppo
\newcommand{\RM}{Riccardo Montagnin}
\newcommand{\MT}{Manuel Turetta}
\newcommand{\FB}{Francesco Bazzerla}
\newcommand{\SL}{Stefano Lia}
\newcommand{\LD}{Luca Dario}
\newcommand{\DC}{Diego Cavestro}
\newcommand{\ND}{Nicolò Dovico}

%Ruoli
\newcommand{\Pm}{Project Manager}
\newcommand{\Am}{Amministratore}
\newcommand{\AmP}{Amministratori}
\newcommand{\An}{Analista}
\newcommand{\AnP}{Analisti}
\newcommand{\Dev}{Sviluppatore}
\newcommand{\DevP}{Sviluppatori}
\newcommand{\Ver}{Verificatore}
\newcommand{\VerP}{Verificatori}
\newcommand{\Progr}{Programmatore}
\newcommand{\ProgrP}{Programmatori}
\newcommand{\Prog}{Progettista}
\newcommand{\ProgP}{Progettisti}



%Firme
\newcommand{\RMFirma}{\myincludegraphics[scale = 0.5]{../../../Template/Firme/RM.png}}
\newcommand{\MTFirma}{\myincludegraphics[scale = 0.5]{../../../Template/Firme/MT.png}}
\newcommand{\FBFirma}{\myincludegraphics[scale = 0.5]{../../../Template/Firme/FB.png}}
\newcommand{\SLFirma}{\myincludegraphics[scale = 0.5]{../../../Template/Firme/SL.png}}
\newcommand{\LDFirma}{\myincludegraphics[scale = 0.5]{../../../Template/Firme/LD.png}}
\newcommand{\DCFirma}{\myincludegraphics[scale = 0.5]{../../../Template/Firme/DC.png}}
\newcommand{\NDFirma}{\myincludegraphics[scale = 0.5]{../../../Template/Firme/ND.png}}

%Professori e proponente
\newcommand{\TV}{Prof. Tullio Vardanega}
\newcommand{\RC}{Prof. Riccardo Cardin}
\newcommand{\RB}{Red Babel}
\newcommand{\proponente}{Red Babel}

%Documenti
\newcommand{\Gl}{Glossario}
\newcommand{\glossario}{\textit{\Gl\_v.2.0.0.pdf}}
\newcommand{\AdR}{Analisi dei Requisiti}
\newcommand{\analisiDeiRequisiti}{\textit{\AdR\_v.2.0.0.pdf}}
\newcommand{\AdRvDue}{AnalisiDeiRequisiti}
\newcommand{\NdP}{Norme di Progetto}
\newcommand{\normeDiProgetto}{\textit{\NdP\_v.2.0.0.pdf}}
\newcommand{\PdP}{Piano di Progetto}
\newcommand{\pianoDiProgetto}{\textit{\PdP\_v.2.0.0.pdf}}
\newcommand{\SdF}{Studio di Fattibilità}
\newcommand{\studioDiFattibilita}{\textit{\SdF\_v.2.0.0.pdf}}
\newcommand{\PdQ}{Piano di Qualifica}
\newcommand{\pianoDiQualifica}{\textit{\PdQ\_v.2.0.0.pdf}}
\newcommand{\VI}{Verbale Interno}
\newcommand{\VE}{Verbale Esterno}
\newcommand{\ST}{Specifica Tecnica}
\newcommand{\MU}{Manuale Utente}
\newcommand{\DDP}{Definizione di Prodotto}

%Periodo di progetto
\newcommand{\ARM}{Analisi dei Requisiti di Massima}
\newcommand{\ARD}{Analisi dei Requisiti in Dettaglio}
\newcommand{\PA}{Progettazione Architetturale}
\newcommand{\PD}{Progettazione di Dettaglio}
\newcommand{\COD}{Codifica}
\newcommand{\VV}{Verifica e Testing Finale}

%Consegne
\newcommand{\RR}{Revisione dei Requisiti}
\newcommand{\RP}{Revisione di Progettazione}
\newcommand{\RQ}{Revisione di Qualifica}
\newcommand{\RA}{Revisione di Accettazione}


%Formattazione
\newcommand{\termine}[1]{\textit{#1}\small{$_G$}}
\newcommand{\link}[1]{\href{#1}{\textcolor{blue}{\texttt{#1}}}} 

% Testi ricorrenti
\newcommand{\scopoProdotto}{L'obiettivo di questo progetto è la realizzazione di un \termine{SDK} che permetta la creazione di bolle interattive, le quali, successivamente, verranno utilizzate all'interno dell'applicazione di messaggistica istantanea open source \termine{Rocket.chat}. \\
Dopo la realizzazione di tale \termine{SDK}, è proposto lo sviluppo di un'applicazione in grado di sfruttare l'\termine{SDK} per implementare un uso originale. L'applicazione scelta dal \termine{team} consiste nella bolla lista-spesa e nei suoi vari utilizzi all'interno della piattaforma \termine{Rocket.chat}.
}
\newcommand{\descrizioneGlossario}{Al fine di mantenere questo documento compatto e di facile lettura è stato realizzato un glossario esterno contenente tutte le definizioni dei termini che più comunemente verranno presentati al lettore.  
Tale glossario si ritrova all'interno del file \glossario, e contiene tutti e soli i termini che vengono marcati con una \textit{G} a pedice.
}
\newcommand{\riferimentiNormativi}{
	\begin{itemize}
		\item \textbf{Norme di Progetto}: \normeDiProgetto
		\item \textbf{\termine{Capitolato} d'appalto C5: Monolith - An Interactive bubble provider} \\
			  \link{http://www.math.unipd.it/~tullio/IS-1/2016/Progetto/C5.pdf}
	\end{itemize}
}

% Comandi per generare l'intro
\newcommand{\documento}{\PdQ}
\newcommand{\versione}{1.0.0}
\newcommand{\redatori}{\RM\\ & \MT\\ & \FB}
\newcommand{\revisori}{\DC\\ & \RM}
\newcommand{\approvazione}{\LD}
\newcommand{\statoapprovazione}{Approvato}
% Quando il documento sarà approvato, inserire all'interno del comando seguente la data nel formato GG mese AAAA dove GG è il giorno a due cifre, mese è il mese scritto per esteso con la prima lettera minuscola, e AAAA è l'anno a quattro cifre
\newcommand{\dataApprovazione}{}
\newcommand{\uso}{Esterno}
\newcommand{\destinatari}{\TV\\ & \RC\\ & \RB}

\newcommand{\sommario}{Questo documento si prefigge di regolamentare le operazioni di verifica del gruppo \gruppo\ necessarie ad assicurare i requisiti qualitativi per il progetto \progetto.
}

\newcommand{\modifiche}{
3.0.0 & Approvazione del documento - Creare nuova versione del documento & \SL & \Pm & 07/05/2017 \\\midrule
2.1.0 & Verifica documento - Correzione errori & \LD & \Ver & 04/05/2017 \\\midrule
2.0.4 & Aggiunta la nota riguardante la forma del \MU\ di \progettoShort\ - In seguito a quanto deciso in data 03 maggio 2017 e riportato nell'apposito verbale & \ND & \Am & 04/05/2017 \\\midrule
	2.0.3 & Modificata la sezione delle metriche per la codifica - Aggiungere regole più precise per facilitare la stesura del codice & \ND & \Am & 28/03/2017 \\\midrule
2.0.2 & Aggiunte sezioni mancanti - Migliorare la profondità del documento come segnalato nella correzione in seguito alla \RP & \SL & \Am & 27/03/2017 \\
\midrule
2.0.1 & Riorganizzata la sezione degli strumenti - Aggiungere chiarezza su quale strumento sia usato per quale attività & \SL & \Am & 26/03/2017 \\
\midrule
2.0.0 & Approvazione del documento - Creare nuova versione del documento & \DC & \Pm & 04/03/2017 \\
\midrule 
1.1.0 & Verifica del documento - Correzioni errori & \LD & \Ver & 04/03/2017 \\
\midrule 
	1.0.2 & Aggiunta metriche - Aggiunta profondità come segnalato nella correzione in seguito alla \RR & \ND & \Am & 28/02/2017 \\
	\midrule
	1.0.1 & Aggiornamenti sezioni 3 e 4 - Aggiunta ampiezza come segnalato nella correzione in seguito alla \RR & \ND & \Am & 27/02/2017 \\
	\midrule
	1.0.0 & Approvazione - Creare la prima versione del documento & \SL & \Pm & 04/01/2016 \\
	\midrule
	0.4.0 & Verifica sottosezione 4.4 - Correzione errori & \RM & \Ver & 29/12/2016 \\\midrule
	0.3.1 & Stesura sottosezione 4.4 - Aggiunte norme sul processo di formazione & \DC & \Ver & 29/12/2016 \\\midrule
	0.3.0 & Verifica sezione 4 - Correzione errori & \RM &\Ver & 27/12/2016 \\\midrule
	0.2.0 & Verifica sezione 3 - Correzione errori & \RM & \Ver & 24/12/2016 \\\midrule
	0.1.0  & Verifica sezioni 1 e 2 - Correzioni errori & \RM & \Ver & 23/12/2016\\\midrule
    0.0.5 & Stesura sezione 4 - stabilire norme per il coordinamento e pianificazione & \DC & \Am & 21/12/2016 \\\midrule
    0.0.4 & Stesura sezione 2 - Definizione dei processi primari per il progetto & \LD & \Am & 20/12/2017 \\\midrule
    0.0.3 & Stesura sezione 1 e modifica del template - Introduzione al documento & \FB & \Am & 18/12/2016 \\\midrule
    0.0.2 & Stesura sezione 3 - Definizione dei processi di supporto  & \ND & \Am & 17/12/2016 \\\midrule
    0.0.1 & Creazione del template - Inizio documento & \SL & \Am & 15/12/2016 \\\midrule
}


\begin{document}

% Questo file contiene il layout della prima pagina
\pagenumbering{gobble}

\title{\includegraphics[width=8cm, keepaspectratio=true]{../../../Template/Logo/Logo.png} \\
	\documento \\
	Versione \versione
}
\date{\dataApprovazione}

\maketitle

\begin{center}

\begin{tabular}{ r | l }
  \textbf{Ruolo} & \textbf{Componente} \\
  Redazione & \redatori \\
  Revisione & \revisori \\
  Approvazione & \approvazione \\
  \\
  Stato & \statoapprovazione \\
  Uso & \uso \\
  Destinatari & \destinatari
\end{tabular}
\end{center}

\begin{center}
\textbf{Sommario\\}
\sommario \\
\vspace{1.5cm}\email
\end{center}

\clearpage

\pagenumbering{arabic}
%Questo file si occupa di generare la tabella delle modifiche
\pagenumbering{Roman}

\begin{center}
    \Large{\textbf{Registro delle modifiche}}
    	\\\vspace{0.5cm}
    	\normalsize
    \begin{tabularx}{\textwidth}{cXXcc}
        \textbf{Versione} & \textbf{Modifica - Motivazione} & \textbf{Autore} & \textbf{Ruolo} & \textbf{Data} \\\toprule
        \modifiche
    \end{tabularx}
\end{center}

\newpage



\tableofcontents

\newpage

\pagenumbering{arabic}



% Sezioni
\section{Introduzione}
\subsection{Scopo del documento}
Questo documento vuole definire le strategie che il \termine{team} ha deciso di adottare per perseguire gli obiettivi di qualità di processo e di prodotto ricercati. A tal fine è necessaria una costante attività di verifica e validazione del lavoro svolto in modo da poter rilevare e correggere le anomalie che potrebbero nascere.

\subsection{Scopo del prodotto}
\scopoProdotto

\subsection{Glossario}
\descrizioneGlossario

\subsection{Riferimenti}
\subsubsection{Normativi}
\riferimentiNormativi

\subsubsection{Informativi}
\begin{itemize}
	\item \textbf{\AdR}: \analisiDeiRequisiti;
	\item \textbf{\PdP}: \pianoDiProgetto;
	\item \textbf{\textit{Slide} dell'insegnamento di Ingegneria del Software}: \\
		  \link{http://www.math.unipd.it/~tullio/IS-1/2016/}
	\item \textbf{\textit{Standard} ISO/IEC 9126}: Product quality \\
	 	  \link{https://en.wikipedia.org/wiki/ISO/IEC\_9126}
	\item \textbf{\textit{Standard} tecnici ISO/IEC 15504}: Software process assessment \\
		  \link{https://en.wikipedia.org/wiki/ISO/IEC\_15504}
	\item \textbf{Ciclo di Deming (\termine{PDCA})}: Miglioramento dei processi \\
		  \link{https://en.wikipedia.org/wiki/PDCA}
\end{itemize}

\newpage
\section{Visione Generale della Strategia di Gestione della Qualità}
\subsection{Obiettivi di Qualità}
\subsubsection{Modello per la Qualità di Processo}
La qualità del prodotto è conseguenza anche della qualità dei processi che lo definiscono. Il \termine{gruppo} però, dopo una attenta valutazione, ha constatato che, essendo una qualità verificabile soltanto a lungo termine, non vi è il tempo materiale per assicurare una qualità di processo definita dall'\termine{ISO}. Il \termine{gruppo}, consapevole della scelta, cercherà comunque di prendere spunto, il più possibile, dallo standard ISO/IEC 15504.

\subsubsection{Standard per la Qualità di Prodotto}
\gruppo\ si impegna a seguire lo standard ISO/IEC 9126 redatto con lo scopo di descrivere obiettivi qualitativi e delineare delle metriche capaci di misurare il raggiungimento di tali obiettivi.

\subsubsection{Soluzioni attuate per il controllo della Qualità di Processo}

L'attuazione del metodo di gestione \termine{PDCA} aiuterà un maggior avvicinamento allo standard ISO/IEC 15504 ed assicurerà, quindi, una maggior qualità di processo. Con il ciclo \termine{PDCA} è possibile infatti garantire un miglioramento continuo dei processi, inclusa la verifica, ed un utilizzo ottimale delle risorse, ottenendo di conseguenza il miglioramento dei prodotti risultanti.
Per avere controllo sulla qualità è necessario che: 
\begin{itemize}
\item
I processi siano pianificati nel dettaglio;
\item
Nella pianificazione siano ripartite in modo chiaro le risorse;
\item
I processi vengano costantemente monitorati.
\end{itemize}

L'attuazione di tali punti è descritta dettagliatamente nel \PdP.
La qualità dei processi viene inoltre monitorata mediante l'analisi costante della qualità del prodotto e quantificata utilizzando le varie metriche che saranno descritte in seguito all'interno di questo documento.

\subsubsection{Soluzioni attuate per il controllo della Qualità di Prodotto}

Al fine di garantire un controllo sistematico della qualità del prodotto, il \termine{gruppo} seguirà le seguenti linee guida:
\begin{itemize}
\item
\textbf{\termine{Quality Assurance}}: insieme di attività realizzate per garantire il raggiungimento degli obiettivi di qualità. Prevede l'attuazione di tecniche di analisi statica e dinamica;
\item
\textbf{\termine{Verifica}}: processo che determina se il lavoro di un determinato periodo è consistente, completo e corretto. La verifica andrà eseguita costantemente durante l'intera durata del progetto;
\item
\textbf{\termine{Validazione}}: conferma in modo oggettivo che il sistema soddisfi correttamente i requisiti. \\ Anch'essa come la \termine{verifica} andrà eseguita costantemente durante l'intera durata del progetto.
\end{itemize}

\subsection{Scadenze Temporali}
Al fine di perseguire l'obbiettivo di rispettare le scadenze fissate nel Piano di Progetto è necessario che l'attività di verifica della documentazione sia sistematica e ben organizzata. Solo così, infatti, l'individuazione e la correzione di eventuali errori avverrà il prima possibile, impedendo la compromissione dell'intero progetto. \\
Ogni attività di redazione dei documenti e di codifica dovrà essere preceduta da uno studio preliminare sulla struttura e sui contenuti degli stessi, con lo scopo di ridurre la possibilità di commettere imprecisioni di natura concettuale e tecnica.

\subsubsection{Responsabilità}

Per ottenere un maggior livello di efficacia ed efficienza nell'attività di verifica verranno attribuite delle responsabilità a specifici ruoli di progetto.
La responsabilità, per l'attività di \termine{verifica} e \termine{validazione}, sarà a carico dei membri del \termine{gruppo} che al momento di eseguire tale attività saranno in carica dei ruoli di \Pm\ e dei \VerP.

\newpage
\section{Strategia di Gestione della Qualità in Dettaglio}
\subsection{Risorse}
Per raggiungere gli obiettivi qualitativi prefissati è necessario, oltre alle risorse umane, utilizzare la potenza e l'affidabilità delle risorse tecnologiche. Infatti, per agevolare il lavoro dei \VerP, verranno impiegati numerosi strumenti automatici che eseguiranno controlli sistematici sui prodotti generati. \\
Le risorse tecniche e tecnologiche consistono in tutti quegli strumenti software e hardware che il \gruppo\ intende utilizzare per le attività di verifica su processi e prodotti; per la descrizione di tali strumenti e delle tecniche utilizzate dal gruppo si rimanda al documento \NdP.

\subsection{Misure e metriche}
Allo scopo di rendere quantificabile il processo di verifica verranno adottate delle misure basate su metriche stabilite a priori. Le metriche incerte, qualora ve ne fossero, verranno migliorate in modo incrementale. \\
Le varie misure che verranno rilevate saranno analizzate confrontandole con due categorie di misurazione:
\begin{itemize}
\item
\textbf{Accettazione}: intervallo di valori minimi entro i quali il prodotto sarà accettato;
\item
\textbf{Ottimale}: intervallo di valori ottimali entro i quali il prodotto risulta soddisfare pienamente, o quasi, i requisiti richiesti. I valori all'interno di tali intervalli devono essere quindi intesi come consigliati ma non vincolanti.
\end{itemize}

\subsection{Strumenti}
Per aiutare la verifica delle metriche descritte nelle \NdP\ verranno utilizzati degli strumenti che, in maniera automatica, svolgeranno dei test e forniranno un resoconto del risultato. Tali strumenti vengono descritti già nelle \NdP\ per cui in questo documento verranno solamente elencati.

\begin{itemize}
\item \textbf{\termine{Jenkins}}
\item \textbf{\termine{SonarQube}}
\item \textbf{\termine{Mocha}:} 
Quest'ultimo affiancato da:
\begin{itemize}
\item \textbf{\termine{Chai}}
\item \textbf{\termine{Sinon}}
\end{itemize}
\end{itemize}

%----Non so se Nico ha copiato questa parte sulle Norme.----
%
%\subsubsection{Jenkins}
%\termine{Jenkins} è un software che permetterà di automatizzare una serie di operazioni che consentiranno al \termine{team} di monitorare la qualità del software in modo continuo ogni volta che verrà inviata una modifica alla \termine{repository} che conterrà il codice sorgente. \\
%Per ottenere questi risultati questo software si appoggerà agli altri strumenti descritti qua di seguito.
%
%\subsubsection{SonarQube}
%\termine{SonarQube} è un software che permette di eseguire analisi statica sul codice che verrà inviato alla \termine{repository} fornendo un feedback al programmatore per fargli sapere cosa dovrà migliorare. Inoltre permetterà di controllare che le regole che il \termine{team} si è imposto per la qualità del codice sorgente vengano rispettate.
%
%\subsubsection{Mocha}
%\termine{Mocha} è un \termine{framework} che permette di eseguire test di unità e di integrazione andando a creare dei \termine{mock} che simulino il comportamento delle varie parti, con lo scopo di testare ogni dettaglio in modo indipendente dal resto. Ad esso si affiancheranno anche:
%\begin{itemize}
%\item \textbf{\termine{Chai}}: libreria per formulare asserzioni sull'output atteso;
%\item \textbf{\termine{Sinon}}: libreria per simulare delle risposte di richieste a server remoti.
%\end{itemize}
\section{Capability Maturity Model}
Il \termine{Capability Maturity Model} (\textit{CMM}) è un modello di riferimento costituito
da pratiche consolidate in una disciplina specifica e viene utilizzato per stabilire la
capacità di una organizzazione o gruppo di operare in quella disciplina.

\subsection{I modelli CMM}
Esistono vari modelli di CMM che differiscono per:

\begin{itemize}
	\item Disciplina (software, sistemi, acquisizione, etc.);
	\item Struttura;
	\item Come viene definito il concetto di \textit{Maturity}, ossia il percorso di miglioramento del processo;
	\item Come viene definito il concetto di \textit{Capability}, ossia l'istituzionalizzazione.
\end{itemize}
Dal 1991 ad oggi sono stati sviluppati dei CMM per molte discipline; alcuni dei
più utilizzati includono modelli per l'ingegneria dei sistemi, l'ingegneria del software, l'acquisizione software, la gestione della forza-lavoro e dello sviluppo ed infine lo sviluppo
integrato di prodotti e processi. \\
L'uso di modelli differenti è risultato problematico a causa della confusione causata dall'utilizzo contemporaneo di più modelli stessi, e dei costi di integrazione di più di un modello in un unico programma coordinato di miglioramento.

\subsection{CMMI}
Per risolvere la confusione che veniva a crearsi con l'utilizzo di più modelli è nato
il progetto di \textit{Integrazione dei CMM} (\termine{Capability Maturity Model Integration -- CMMI}). \\
Nel modello CMMI si usano i livelli per descrivere un percorso di evoluzione raccomandato per un'organizzazione che voglia migliorare i processi che usa per sviluppare e mantenere i propri prodotti e servizi. I livelli presenti sono i seguenti:
\begin{itemize}
	\item \textbf{Livello 0}: Incomplete
	\item \textbf{Livello 1}: Performed
	\item \textbf{Livello 2}: Managed
	\item \textbf{Livello 3}: Defined
	\item \textbf{Livello 4}: Quantitatively Managed
	\item \textbf{Livello 5}: Optimizing
\end{itemize}

\subsubsection{Livello 0: Incomplete}
Un processo è incompleto quando non è stato eseguito oppure è stato eseguito solo parzialmente. Ciò significa che uno o più obiettivi specifici dell'area di processo non sono stati soddisfatti e/o non esistono obiettivi generici per questo livello in quanto non c'è ragione di istituzionalizzare un processo parzialmente eseguito.

\subsubsection{Livello 1: Performed}
Affinché un'area di processo sia al 1$^{\circ}$ livello di \textit{capability} bisogna mettere in pratica
le attività di base richieste per iniziare a lavorare su quella determinata area. \\
Un processo eseguito è tale quando soddisfa gli obiettivi specifici dell'area di processo.
Questo livello supporta e permette le attività necessarie per produrre i \termine{work product} e benché il processo eseguito abbia per risultato dei miglioramenti importanti, questi
potrebbero essere persi nel tempo se non vengono istituzionalizzati. L'istituzionalizzazione, ovvero il soddisfacimento delle pratiche generiche dei livelli di \textit{capability} dal 2 al 5, aiuta inoltre ad assicurare che i miglioramenti siano mantenuti.

\subsubsection{Livello 2: Managed}
Un processo viene definito gestito se viene eseguito e possiede le infrastrutture di base per
essere supportato. Ciò può verificarsi quando esso è organizzato ed eseguito in accordo ad alcune regole, impiega persone competenti che abbiano adeguate conoscenze per produrre dei risultati controllati, coinvolge gli \termine{stakeholder} più importanti ed è monitorato, controllato e revisionato. \\
Questo livello aiuta ad assicurare che le pratiche esistenti siano mantenute nel tempo.

\subsubsection{Livello 3: Defined}
Un processo definito è un processo gestito che è stato adattato in accordo con le linee guida di \termine{tailoring} dei processi standard dell'organizzazione. Esso contribuisce inoltre alla realizzazione dei \termine{work product}, alle misure ed ad altre informazioni di miglioramento. \\
I processi vengono gestiti in maniera preventiva attraverso la comprensione delle relazioni tra le attività di processo, le misure dettagliate, i \termine{work product} ed i servizi.

\subsubsection{Livello 4: Quantitatively Managed}
In questo livello il processo è un processo definito e controllato usando tecniche quantitative o statistiche. Gli obiettivi quantitativi per la qualità del processo sono stabiliti ed usati come criteri di gestione dei processi stessi. Infine, la performance del processo e della qualità è intesa in termini statistici e viene gestita attraverso tutta la vita del processo.

\subsubsection{Livello 5: Optimizing}
In questo caso si ha un livello di \textit{capability} pari a quello del livello precedente, con l'unica differenza che è stato migliorato sulla base di una comprensione delle più comuni cause di variazione del processo. Il punto focale a questo livello è un miglioramento continuo del range di esecuzione del processo attraverso miglioramenti innovativi ed incrementali.

\subsection{Scelte del team}
A causa della ridotta quantità di tempo messa a disposizione del \termine{team} per la realizzazione del progetto, è stato deciso unanimemente che non verrà eseguita alcuna attività con lo scopo di migliorare la qualità dei processi ma, allo stesso tempo, tutti i membri del \termine{gruppo} svolgeranno i propri compiti con l'obiettivo di provare a soddisfare sempre il livello 2 del \textit{CMMI}, senza però la garanzia di un totale raggiungimento di questo obiettivo. 

\newpage
\section{Standard ISO/IEC9126}
Lo standard ISO/IEC 9126 è stato redatto con lo scopo di descrivere quali sono gli obiettivi qualitativi che un prodotto software deve soddisfare. Questi vengono suddivisi in 3 aree tematiche diverse:
\begin{itemize}
	\item\textbf{Qualità esterna}: rappresenta la qualità del software nel momento in cui esso viene eseguito e testato;
	\item\textbf{Qualità interna}: rappresenta la qualità del software per quanto riguarda le sue caratteristiche implementative, durante le fasi di progettazione e codifica;
	\item\textbf{Qualità in uso}: rappresenta la qualità del software dal punto di vista del cliente che lo sta utilizzando.
\end{itemize}

Lo standard delinea sei macro-obbiettivi qualitativi, i quali sono suddivisi a loro volta in sotto caratteristiche specifiche:
\begin{itemize}
	\item\textbf{Funzionalità}: capacità del prodotto di fornire tutte le funzioni che sono state individuate attraverso l'\AdR:
	\begin{itemize}
		\item\textbf{Adeguatezza}: le funzionalità fornite devono essere conformi rispetto le aspettative;
		\item\textbf{Accuratezza}: il prodotto deve fornire i risultati attesi, con il livello di precisione richiesto;
		\item\textbf{Interoperabilità}: il prodotto deve poter interagire ed operare con uno o più sistemi specifici;
		\item\textbf{Sicurezza}: il prodotto deve proteggere le informazioni e i dati da accessi e modifiche non autorizzati;
		\item\textbf{Conformità di funzionalità}: il prodotto deve aderire a standard, regole e convenzioni inerenti alla funzionalità.
	\end{itemize}

	\item\textbf{Affidabilità}: capacità del prodotto software di svolgere correttamente le sue funzioni durante il suo utilizzo, anche nel caso in cui si presentino situazioni anomale:
	\begin{itemize}
		\item\textbf{Maturità}: il prodotto deve evitare che si verifichino malfunzionamenti o che vengano prodotti risultati non corretti;
		\item\textbf{Tolleranza agli errori}: nel caso in cui si presentino degli errori, dovuti a guasti o ad un uso scorretto dell'applicativo, questi devono essere gestiti in modo da mantenere alto il livello di prestazione;
		\item\textbf{Recuperabilità}: il prodotto deve essere in grado di ristabilire un
adeguato livello di prestazioni e di recuperare i dati rilevanti in seguito a errori o malfunzionamenti;
		\item\textbf{Conformità di affidabilità}: il prodotto deve aderire a standard, regole e convenzioni inerenti all'affidabilità.
	\end{itemize}

	\item\textbf{Usabilità}: capacità del prodotto di essere facilmente comprensibile e attraente in ogni sua parte per qualsiasi utente che lo andrà ad utilizzare:
	\begin{itemize}
		\item\textbf{Comprensibilità}: l'utente deve essere in grado di riconoscerne le funzionalità offerte dal software e deve comprenderne le modalità di utilizzo per riuscire a raggiungere i risultati attesi;
		\item\textbf{Apprendibilità}: deve essere data la possibilità all'utente di imparare ad utilizzare l'applicazione senza troppo impegno;
		\item\textbf{Operabilità}: le funzionalità presenti devono essere coerenti con le aspettative dell'utente;
		\item\textbf{Attrattiva}: il software deve essere piacevole per chi ne fa uso;
		\item\textbf{Conformità di usabilità}: il prodotto deve aderire a standard, regole e convenzioni inerenti all'usabilità.
	\end{itemize}

	\item\textbf{Efficienza}: capacità di eseguire le funzionalità offerte nel minor tempo possibile utilizzando al tempo stesso il minor numero di risorse possibili:
	\begin{itemize}
		\item\textbf{Comportamento rispetto al tempo}: per svolgere le sue funzioni il software deve fornire adeguati tempi di risposta ed elaborazione;
		\item\textbf{Utilizzo delle risorse}: il software quando esegue le sue funzionalità deve utilizzare un appropriato numero e tipo di risorse;
		\item\textbf{Conformità di efficienza}: il prodotto deve aderire a standard, regole e convenzioni inerenti all'efficienza.
	\end{itemize}

	\item\textbf{Manutenibilità}: capacità del prodotto di essere modificato, tramite correzioni, miglioramenti o adattamenti del software a cambiamenti negli ambienti, nei requisiti e nelle specifiche funzionali:
	\begin{itemize}
		\item\textbf{Analizzabilità}: il software deve consentire una rapida identificazione delle possibili cause di errori e malfunzionamenti;
		\item\textbf{Modificabilità}: il prodotto originale deve permettere eventuali cambiamenti in alcune sue parti;
		\item\textbf{Stabilità}: non devono insorgere effetti indesiderati in seguito a modifiche effettuate sul software;
		\item\textbf{Testabilità}: il software deve poter essere facilmente testato per validare le modifiche effettuate;
		\item\textbf{Conformità di manutenibilità}: il prodotto deve aderire a standard, regole e convenzioni inerenti alla manutenibilità.
	\end{itemize}

	\item\textbf{Portabilità}: capacità del software di poter essere utilizzato su diversi ambienti:
	\begin{itemize}
		\item\textbf{Adattabilità}: il prodotto deve adattarsi a tutti quegli ambienti di lavoro nei quali è stato previsto un suo utilizzo, senza dover apportare modifiche allo stesso;
		\item\textbf{Installabilità}: il software deve poter essere installato in determinati ambienti di lavoro;
		\item\textbf{Coesistenza}: il prodotto può coesistere in ambienti comuni
con altri software, condividendo risorse comuni;
		\item\textbf{Sostituibilità}: l'applicativo deve poter sostituire un altro software che ha lo stesso scopo e lavora nel medesimo ambiente;
		\item\textbf{Conformità di portabilità}: il prodotto deve aderire a standard, regole e convenzioni inerenti alla portabilità.
	\end{itemize}
\end{itemize}

\subsection{Scelte del team}
Per garantire una buona qualità di prodotto, il \termine{team} ha individuato dallo standard \textit{ISO/IEC 9126} le qualità che ritiene più importanti nell'arco del ciclo di vita del prodotto e le ha istanziate individuando obiettivi e metriche coerenti con i livelli di qualità perseguiti.

\newpage
\section{PDCA}
Il ciclo \termine{PDCA}, detto anche \termine{Ciclo di Deming}, definisce un metodo di controllo dei processi durante il loro ciclo di vita che consente di migliorarne continuamente la qualità. \\
Tale approccio è suddiviso in 4 fasi:
\begin{itemize}
\item \textbf{Plan}: fase di pianificazione, dove si individuano gli obiettivi e i processi necessari per il raggiungimento dei risultati attesi;
\item \textbf{Do}: fase di attuazione del piano individuato al passo precedente e raccolta di dati sulla qualità ottenuta;
\item \textbf{Check}: fase di verifica, dove si confrontano i risultati ottenuti (fase di Do) ed i risultati attesi (fase di Plan);
\item \textbf{Act}: fase in cui si determinano le cause delle differenze fra risultati ottenuti e risultati attesi, per decidere dove attuare eventuali azioni correttive per avere un effettivo miglioramento della qualità.
\end{itemize}

A causa della mancanza di tempo e di risorse è stato deciso unanimemente da tutti i componenti del \termine{team} che non verrà istanziato alcun processo di automiglioramento della qualità dei processi ma, allo stesso tempo, ogni componente si impegnerà a svolgere tutti i compiti lui assegnati nel miglior modo possibile.

\newpage
\section{Test}
Il \termine{team} \gruppo\ al fine di creare del software che sia di qualità ha strutturato dei test atti a verificare
che il software prodotto rispecchi le funzionalità a fronte di risultati attesi.
Tutte le attività di testing prodotte devono poter essere ripetibili e devono essere deterministiche, al fine di poter
fornire delle informazioni utili a intraprendere azioni di correzione nel caso in cui i risultati ottenuti siano diversi da quelli attesi. \\
Per avere un tracciamento dei test prodotti e dei risultati ottenuti si è scelto di classificare il tutto producendo dei
\termine{log} che siano di facile consultazione e che possano fornire una precisa indicazione di quelli che sono stati
gli output di queste attività di verifica, eventuali errori o eventuali risultati che siano non coerenti con quanto in precedenza fissato.

\subsection{Tipi di test}
Sono stati individuati quattro livelli di testing:
\begin{itemize}
\item \textbf{Test di unità [TU]:} con questa tipologia di test si cerca di verificare la più piccola parte di lavoro prodotta da un programmatore. Questo si traduce tendenzialmente a verificare i metodi e le funzioni scritte;
\item \textbf{Test di integrazione [TI]:} con questa tipologia di test si cerca di verificare le componenti di sistema. Più precisamente, l'obiettivo è quello di testare il funzionamento dei vari package prodotti, sia singolarmente che nel loro insieme;
\item \textbf{Test di sistema [TS]:} con questa tipologia di test si cerca di verificare che il comportamento e il funzionamento dell'architettura siano corretti;
\item \textbf{Test di validazione [TV]:} con questa tipologia di test si vuole verificare che il lavoro prodotto soddisfi quanto richiesto dal proponente.
\end{itemize}

% INSERIMENTO DEI TEST
\subsubsection{Test di Unità}
Con questa tipologia di test si vuole verificare il corretto funzionamento delle unità individuate durante la definizione dell'architettura.
I test di unità saranno descritti nel modo seguente:
\begin{center}
\textbf{TU}[IdTest]
\end{center}
dove:
\begin{itemize}
\item IdTest rappresenta il codice identificativo crescente dell'unità considerata.
\end{itemize}

\begin{center}
	\begin{longtable}{|c|>{\centering}m{10cm}|c|c|}\hline
		Id & Descrizione & Stato & Esito\\ \hline
		TU1 & Verificare che in un widget di tipo testo venga impostata in modo corretto la grandezza del font del testo & \textcolor{Green}{Implementato} & \textcolor{Green}{Superato} \\ \hline
		TU2 & Verificare che all'instanziazione di un widget di tipo testo la grandezza del testo sia di default a 1em & \textcolor{Green}{Implementato} & \textcolor{Green}{Superato} \\ \hline
		TU3 & Verificare che all'interno di un widget di tipo testo venga impostata in corsivo la parte di testo richiesta & \textcolor{Green}{Implementato} & \textcolor{Green}{Superato} \\ \hline
		TU4 & Verificare che all'interno di un widget di tipo testo venga inserito un link cliccabile nel modo corretto & \textcolor{Green}{Implementato} & \textcolor{Green}{Superato} \\ \hline
		TU5 & Verificare che all'interno di un widget di tipo testo i link cliccabili abbiano il colore corretto & \textcolor{Green}{Implementato} & \textcolor{Green}{Superato} \\ \hline
		TU6 & Verificare che all'interno di un widget di tipo testo venga impostata in grassetto la parte di testo richiesta & \textcolor{Green}{Implementato} & \textcolor{Green}{Superato} \\ \hline
		TU7 & Verificare che all'interno di un widget di tipo testo venga impostato il colore del testo in modo corretto & \textcolor{Green}{Implementato} & \textcolor{Green}{Superato} \\ \hline
		TU8 & Verificare che all'instanziazione di un widget di tipo testo il colore del testo sia di default a nero & \textcolor{Green}{Implementato} & \textcolor{Green}{Superato} \\ \hline
		TU9 & Verificare che in un widget di tipo immagine l'immagine venga agiunta in modo corretto & \textcolor{Green}{Implementato} & \textcolor{Green}{Superato} \\ \hline
		TU10 & Verificare che in un widget di tipo immagine venga mostrato il messaggio di errore appropriato nel caso si inserisca un immagina non valida & \textcolor{Green}{Implementato} & \textcolor{Green}{Superato} \\ \hline
		TU11 & Verificare che un widget di tipo immagine venga impostata la dimensione del immagine da visualizzare in modo corretto & \textcolor{Green}{Implementato} & \textcolor{Green}{Superato} \\ \hline
		TU12 & Verificare che all'instanziazione di un widget di tipo immagine venga impostata la larghezza e l'altezza di default & \textcolor{Green}{Implementato} & \textcolor{Green}{Superato} \\ \hline
		TU13 & Verificare un widget di tipo immagine venga visualizzato con le dimensioni corrette & \textcolor{Green}{Implementato} & \textcolor{Green}{Superato} \\ \hline
		TU14 & Verificare che in un widget di tipo bottone venga impostato il testo in modo corretto & \textcolor{Green}{Implementato} & \textcolor{Green}{Superato} \\ \hline
		TU15 & Verificare che venga mostrarto in modo corretto il messaggio di errore se viene inserita una sequenza di caratteri non valida in un widget di tipo bottone & \textcolor{Green}{Implementato} & \textcolor{Green}{Superato} \\ \hline
		TU16 & Verificare che un widget di tipo bottone venga impostata la dimensione in modo corretto & \textcolor{Green}{Implementato} & \textcolor{Green}{Superato} \\ \hline
		TU17 & Verificare che all'instanziazione di un widget di tipo bottone venga impostata la larghezza e l'altezza di default & \textcolor{Green}{Implementato} & \textcolor{Green}{Superato} \\ \hline
		TU18 & Verificare che in un widget di tipo bottone venga impostato il colore di sfondo in modo corretto & \textcolor{Green}{Implementato} & \textcolor{Green}{Superato} \\ \hline
		TU19 & Verificare che all'instanziazione di un widget di tipo bottone venga impostato il colore di sfondo di default & \textcolor{Green}{Implementato} & \textcolor{Green}{Superato} \\ \hline
		TU20 & Verificare che in un widget di tipo button vengano eseguiti gli eventi associati alla pressione & \textcolor{Green}{Implementato} & \textcolor{Green}{Superato} \\ \hline
		TU21 & Verificare che in un widget di tipo button vengano eseguiti gli eventi associati alla pressione pressione prolungata & \textcolor{Green}{Implementato} & \textcolor{Green}{Superato} \\ \hline
		TU22 & Verificare che in un widget di tipo button venga impostato in modo corretto la soglia per determinare che una pressione da parte del utente è di tipo: "pressione prolungata" & \textcolor{Green}{Implementato} & \textcolor{Green}{Superato} \\ \hline
		TU23 & Verificare che un un widget di tipo checklistitem venga cambiato in modo corretto la stato da "cliccato" a "non cliccato" e viceversa & \textcolor{Green}{Implementato} & \textcolor{Green}{Superato} \\ \hline
		TU24 & Verificare che all'instanziazione di un widget di tipo checklistitem venga impostato lo stato a "non cliccato" & \textcolor{Green}{Implementato} & \textcolor{Green}{Superato} \\ \hline
		TU25 & Verificare che un widget di tipo checklistitem venga visualizzato con lo stile di spunta e colore corretto & \textcolor{Green}{Implementato} & \textcolor{Green}{Superato} \\ \hline
		TU26 & Verificare che in un widget di tipo checklistitem vengano associati gli eventi di pressione nel modo corretto & \textcolor{Green}{Implementato} & \textcolor{Green}{Superato} \\ \hline
		TU27 & Verificare che in un widget di tipo checklistitem vengano associati gli eventi di pressione prolungata nel modo corretto & \textcolor{Green}{Implementato} & \textcolor{Green}{Superato} \\ \hline
		TU28 & Verificare che in un widget di tipo checklistitem vengano eseguiti gli eventi associati alla pressione & \textcolor{Green}{Implementato} & \textcolor{Green}{Superato} \\ \hline
		TU29 & Verificare che in un widget di tipo checklistitem vengano eseguiti gli eventi associati alla pressione pressione prolungata & \textcolor{Green}{Implementato} & \textcolor{Green}{Superato} \\ \hline
		TU30 & Verificare che in un widget di tipo checklistitem venga impostato in modo corretto la soglia per determinare che una pressione da parte del utente è di tipo: "pressione prolungata" & \textcolor{Green}{Implementato} & \textcolor{Green}{Superato} \\ \hline
		TU31 & Verificare che in un widget di tipo list venga impostato correttamente il cambiamento dell'indicatore a 'pallino' & \textcolor{Green}{Implementato} & \textcolor{Green}{Superato} \\ \hline
		TU32 & Verificare che in un widget di tipo list venga impostato correttamente il cambiamento dell'indicatore a 'trattino' & \textcolor{Green}{Implementato} & \textcolor{Green}{Superato} \\ \hline
		TU33 & Verificare che in un widget di tipo list venga impostato correttamente il cambiamento dell'indicatore ad elenco numerato & \textcolor{Green}{Implementato} & \textcolor{Green}{Superato} \\ \hline
		TU34 & Verificare che in un widget di tipo list che vengano aggiunti correttamente tutti gli elementi & \textcolor{Green}{Implementato} & \textcolor{Green}{Superato} \\ \hline
		TU35 & Verificare che in un layout verticale sia possibile aggiungere un elemento & \textcolor{Green}{Implementato} & \textcolor{Green}{Superato} \\ \hline
		TU36 & Verificare che in un layout verticale non sia possibile aggiungere se stesso & \textcolor{Green}{Implementato} & \textcolor{Green}{Superato} \\ \hline
		TU37 & Verificare che in un layout orizzontale sia possibile aggiungere un elemento & \textcolor{Green}{Implementato} & \textcolor{Green}{Superato} \\ \hline
		TU38 & Verificare che in un layout orizzontale non sia possibile aggiungere se stesso & \textcolor{Green}{Implementato} & \textcolor{Green}{Superato} \\ \hline
		TU39 & Verificare che il nome della lista della spesa venga impostato nel modo corretto & \textcolor{Green}{Implementato} & \textcolor{Green}{Superato} \\ \hline
		TU40 & Verificare che all'instanziazione di una lista della spesa venga impostato il valore di default per il nome della lista & \textcolor{Green}{Implementato} & \textcolor{Green}{Superato} \\ \hline
		TU41 & Verificare che l'immagine della lista della spesa venga impostata in modo corretto & \textcolor{Orange}{Non implementato} & \textcolor{Orange}{Non superato} \\ \hline
		TU42 & Verificare che nel caso di un errore durante l'aggiunta di un prodotto alla lista della spesa venga mostrato un mesaggio di errore & \textcolor{Green}{Implementato} & \textcolor{Green}{Superato} \\ \hline
		TU43 & Verificare che l'immagine del prodotto venga aggiunta in modo corretto & \textcolor{Green}{Implementato} & \textcolor{Green}{Superato} \\ \hline
		TU44 & Verificare che la descrizione del prodotto venga aggiunta in modo corretto & \textcolor{Green}{Implementato} & \textcolor{Green}{Superato} \\ \hline
		TU45 & Verificare che le note del prodotto vengano aggiunte in modo corretto & \textcolor{Orange}{Non implementato} & \textcolor{Orange}{Non superato} \\ \hline
		TU46 & Verificare che la quantità del prodotto venga impostata in modo corretto quando inserita & \textcolor{Green}{Implementato} & \textcolor{Green}{Superato} \\ \hline
		TU47 & Verificare che alla creazione di un nuovo prodotto venga impostata la quantità di default & \textcolor{Green}{Implementato} & \textcolor{Green}{Superato} \\ \hline
		TU48 & Verificare che l'unità di misura relativa alla quantità del prodotto venga impostata in modo corretto & \textcolor{Green}{Implementato} & \textcolor{Green}{Superato} \\ \hline
		TU49 & Verificare che alla creazione di un nuovo prodotto venga impostata l'unità di misura di default & \textcolor{Green}{Implementato} & \textcolor{Green}{Superato} \\ \hline
		TU50 & Verificare che DatabaseSource ritorni la lista che corrisponde all'id inserito & \textcolor{Green}{Implementato} & \textcolor{Green}{Superato} \\ \hline
		TU51 & Verificare che DatabaseSource ritorni l'oggetto che corrisponde all'id inserito & \textcolor{Green}{Implementato} & \textcolor{Green}{Superato} \\ \hline
		TU52 & Verificare che DatabaseSource cancelli la lista che corrisponde all'id inserito & \textcolor{Green}{Implementato} & \textcolor{Green}{Superato} \\ \hline
		TU53 & Verificare che DatabaseSource salvi correttamente la lista nel database & \textcolor{Green}{Implementato} & \textcolor{Green}{Superato} \\ \hline
		TU54 & Verificare che ModifyListUseCase aggiunga correttamente un item alla lista nel database & \textcolor{Green}{Implementato} & \textcolor{Green}{Superato} \\ \hline
		TU55 & Verificare che ModifyListUseCase rimuova correttamente un item alla lista nel database & \textcolor{Green}{Implementato} & \textcolor{Green}{Superato} \\ \hline
		TU56 & Verificare che ModifyListUseCase aggiorni correttamente un item alla lista nel database & \textcolor{Green}{Implementato} & \textcolor{Green}{Superato} \\ \hline
		TU57 & Verificare che DeleteListViewPresenter funzioni correttamente senza lanciare eccezioni & \textcolor{Green}{Implementato} & \textcolor{Green}{Superato} \\ \hline
		TU58 & Verificare che InputListInfoViewPresenter crei correttamente un elemento della lista & \textcolor{Green}{Implementato} & \textcolor{Green}{Superato} \\ \hline
		TU59 & Verificare che ShareWithGroupViewPresenter funzioni correttamente senza lanciare eccezioni & \textcolor{Green}{Implementato} & \textcolor{Green}{Superato} \\ \hline
		TU60 & Verificare che ShareWithContactViewPresenter funzioni correttamente senza lanciare eccezioni & \textcolor{Green}{Implementato} & \textcolor{Green}{Superato} \\ \hline
	\end{longtable}
\end{center}

\subsubsection{Test di Integrazione}
Con questa tipologia di test si vuole determinare il corretto funzionamento delle componenti progettate durante la
definizione dell'architettura ad alto livello.

I test di integrazione saranno descritti nel modo seguente:
\begin{center}
\textbf{TI}[\textit{IdComponente}]
\end{center}
dove:
\begin{itemize}
\item
\textbf{IdComponente} rappresenta il codice identificativo crescente del componente considerato.
\end{itemize}
È stato scelto di utilizzare un approccio top-down nel determinare i test di integrazione.

\begin{center}
	\begin{longtable}{|c|>{\centering}m{10cm}|c|c|}\hline
		Id & Descrizione & Stato & Esito\\ \hline
		TI1 & Si verifica che il metodo ButtonWidget::setText passi correttamente il testo al presenter il quale eventualmente lo formatterà e modificherà correttamente l'HTML e CSS & \textcolor{Green}{Implementato} & \textcolor{Green}{Superato} \\ \hline
		TI2 & Si verifica che il metodo ButtonWidget::setHeight passi correttamente il valore dell'altezza al presenter il quale modificherà correttamente l'HTML e CSS & \textcolor{Green}{Implementato} & \textcolor{Green}{Superato} \\ \hline
		TI3 & Si verifica che il metodo ButtonWidget::setWidth passi correttamente il valore della lareghezza al presenter modificherà correttamente l'HTML e CSS & \textcolor{Green}{Implementato} & \textcolor{Green}{Superato} \\ \hline
		TI4 & Si verifica che il metodo ButtonWidget::setBackgroundColor passi correttamente il colore dello sfondo del bottone al presenter il quale modificherà correttamente l'HTML e CSS & \textcolor{Green}{Implementato} & \textcolor{Green}{Superato} \\ \hline
		TI5 & Si verifica che il metodo ButtonWidget::setOnClickAction passi correttamente al presenter la funzione che verrà eseguita quanto viene premuto il pulsante & \textcolor{Green}{Implementato} & \textcolor{Green}{Superato} \\ \hline
		TI6 & Si verifica che il metodo ButtonWidget::setOnLongClickAction passi correttamente al presenter la funzione che verrà eseguita quanto viene mantenuto premuto il pulsante & \textcolor{Green}{Implementato} & \textcolor{Green}{Superato} \\ \hline
		TI7 & Si verifica che il metodo TextWidget::setText passi correttamente al presenter il testo che verrà poi mostrato a video & \textcolor{Green}{Implementato} & \textcolor{Green}{Superato} \\ \hline
		TI8 & Si verifica che il metodo TextWidget::setTextColor passi correttamente il colore del testo al presenter il quale modificherà correttamente l'HTML e CSS, presentandolo con il colore corrispondente & \textcolor{Green}{Implementato} & \textcolor{Green}{Superato} \\ \hline
		TI9 & Si verifica che il metodo TextWidget::setFormatText passi il valore booleano (che rappresenta la scelta di formattare il testo) al presenter. Esso dovrà, poi, presentare il testo formattato seguendo la sintassi markdown oppure lasciarlo invariato in funzione al valore booleano & \textcolor{Green}{Implementato} & \textcolor{Green}{Superato} \\ \hline
		TI10 & Si verifica che il metodo TextWidget::setUrlHighlightColor passi il valore del colore al presenter il quale dovrà presentare i link evidenziati oppure non evidenziati in funzione al colore scelto & \textcolor{Green}{Implementato} & \textcolor{Green}{Superato} \\ \hline
		TI11 & Si verifica che il metodo TextWidget::setTextSize passi la dimensione del testo del presenter il quale dovrà occuparsi di presentare il testo della dimensione impostata & \textcolor{Green}{Implementato} & \textcolor{Green}{Superato} \\ \hline
		TI12 & Si verifica che il metodo ChecklistItemWidget::setUseSelectionMark imposti la scelta di utilizzare un carattere o un colore come spunta passando correttamente per il \termine{Presenter} & \textcolor{Green}{Implementato} & \textcolor{Green}{Superato} \\ \hline
		TI13 & Si verifica che il metodo ChecklistItemWidget::setSelectionColor passi al presenter il colore dello spunta il quale dovrà presentare la lista con lo stile corretto & \textcolor{Green}{Implementato} & \textcolor{Green}{Superato} \\ \hline
		TI14 & Si verifica che il metodo ChecklistItemWidget::setSelectionCharacter passi al presenter un carattere per la spunta, e il presenter dovrà presentare la lista con lo stile corretto & \textcolor{Green}{Implementato} & \textcolor{Green}{Superato} \\ \hline
		TI15 & Si verifica che il metodo ListWidget::addItem passi al presenter il testo del elemento il quale dovrà presentare l'HTML con il nuovo elemento & \textcolor{Green}{Implementato} & \textcolor{Green}{Superato} \\ \hline
		TI16 & Si verifica che la view passi al presenter il carattere con il quale mostrare gli indicatori della lista. Esso provvederà poi ad impostarli correttamente nell'HTML & \textcolor{Green}{Implementato} & \textcolor{Green}{Superato} \\ \hline
		TI17 & Si verifica che il metodo ImageWidget::setWidth passi correttamente al presenter il valore dell'altezza il quale aggiornerà la view & \textcolor{Green}{Implementato} & \textcolor{Green}{Superato} \\ \hline
		TI18 & Si verifica che il metodo ImageWidget::setHeight passi correttamente al presenter il valore della larghezza il quale aggiornerà la view & \textcolor{Green}{Implementato} & \textcolor{Green}{Superato} \\ \hline
		TI19 & Si verifica che il metodo ImageWidget::setImage passi correttamente al presenter il path del'immagine che dovrà essere presentata nella view & \textcolor{Green}{Implementato} & \textcolor{Green}{Superato} \\ \hline
		TI20 & Si verifica che l'aggiunta di un VerticalLayoutView, creato con degli elementi al suo interno, non causi errori nell'aggiunta di questo alla bolla e che, gli elementi, nel layout appena aggiunto, vengano mostrati uno sotto l'altro & \textcolor{Green}{Implementato} & \textcolor{Green}{Superato} \\ \hline
		TI21 & Si verifica che l'aggiunta di un HorizontalLayoutView, creato con degli elementi al suo interno, non causi errori nell'aggiunta di questi alla bolla e che gli elementi, nel layout appena aggiunto, vengano mostrati uno vicino all'altro & \textcolor{Green}{Implementato} & \textcolor{Green}{Superato} \\ \hline
		TI22 & Si verifica che il ModifyListUseCase esegua le operazioni di aggiunta, rimozione e aggiornamento dei dati trammite DatabaseSource in modo corretto & \textcolor{Green}{Implementato} & \textcolor{Green}{Superato} \\ \hline
		TI23 & Si verifica che ShareWithGroupViewPresenter interagisca in modo corretto con ChatSource & \textcolor{Green}{Implementato} & \textcolor{Green}{Superato} \\ \hline
		TI24 & Si verifica che CreateListViewPresenter interagisca con ChatSource in modo corretto & \textcolor{Orange}{Non implementato} & \textcolor{Orange}{Non superato} \\ \hline
		TI25 & Si verifica che ShareListUseCase interagisca in modo corretto con DatabaseSource & \textcolor{Green}{Implementato} & \textcolor{Green}{Superato} \\ \hline
		TI26 & Si verifica che ManageListUseCase interagisca con DatabaseSource rimuovendo e aggiungendo una lista dal database in modo corretto & \textcolor{Green}{Implementato} & \textcolor{Green}{Superato} \\ \hline
		TI27 & Si verifica che DeleteListViewPresenter interagisca con ShowPopupUseCase mostrando il popup di eliminazione della lista & \textcolor{Green}{Implementato} & \textcolor{Green}{Superato} \\ \hline
		TI28 & Si verifica che ForwardListUseCase interagisca in modo corretto con ChatSource al fine di inoltrare la lista selezionata & \textcolor{Orange}{Non implementato} & \textcolor{Orange}{Non superato} \\ \hline
		TI29 & Si verifica che GetListInfoUseCase esegua i fetch dei dati in modo corretto dal database tramite DatabaseSource & \textcolor{Green}{Implementato} & \textcolor{Green}{Superato} \\ \hline
		TI30 & Si verifica che ChangeListInfoViewPresenter interagisca in modo corretto con ModifyListUseCase & \textcolor{Orange}{Non implementato} & \textcolor{Orange}{Non superato} \\ \hline
		TI31 & Si verifica che ChangeListInfoViewPresenter interagisca in modo corretto con ShowPopupUseCase & \textcolor{Orange}{Non implementato} & \textcolor{Orange}{Non superato} \\ \hline
		TI32 & Si verifica che ShareWithGroupViewImpl interagisca in modo corretto con ShareWithGroupViewPresenter & \textcolor{Green}{Implementato} & \textcolor{Green}{Superato} \\ \hline
		TI33 & Si verifica che ShareWithContactViewImpl interagisca in modo corretto con ShareWithContactViewPresenter & \textcolor{Green}{Implementato} & \textcolor{Green}{Superato} \\ \hline
		TI34 & Si verifica che GetItemInfoUseCase interagisca in modo corretto con DatabaseSource & \textcolor{Green}{Implementato} & \textcolor{Green}{Superato} \\ \hline
		TI35 & Si verifica che ModifyItemPresenter interagisca in modod corretto con ModifyListUseCase & \textcolor{Orange}{Non implementato} & \textcolor{Orange}{Non superato} \\ \hline
		TI36 & Si verifica che ModifyItemPresenter interagisca in modod corretto con ShowPopupUseCase & \textcolor{Orange}{Non implementato} & \textcolor{Orange}{Non superato} \\ \hline
		TI37 & Si verifica che ModifyItemView interagisca in modod corretto con ModifyItemPresenter & \textcolor{Orange}{Non implementato} & \textcolor{Orange}{Non superato} \\ \hline
		TI38 & Si verifica che ModifyListUseCase interagisca in modo corretto con DatabaseSource & \textcolor{Green}{Implementato} & \textcolor{Green}{Superato} \\ \hline
		TI39 & Si verifica che ShareWithContactViewPresenter interagisca in modo corretto con ChatSource & \textcolor{Green}{Implementato} & \textcolor{Green}{Superato} \\ \hline
		TI40 & Si verifica che ShareWithContactViewPresenter interagisca in modo corretto con ShowPopupUseCase & \textcolor{Green}{Implementato} & \textcolor{Green}{Superato} \\ \hline
		TI41 & Si verifica che CreateListViewImpl interagisca in modo corretto con CreateListViewPresenter & \textcolor{Green}{Implementato} & \textcolor{Green}{Superato} \\ \hline
		TI42 & Si verifica che DeleteListViewImpl interagisca in modo corretto con DeleteListViewPresenter & \textcolor{Green}{Implementato} & \textcolor{Green}{Superato} \\ \hline
		TI43 & Si verifica che InputListInfoViewImpl interagisca in modo corretto con InputListInfoViewPresenter & \textcolor{Green}{Implementato} & \textcolor{Green}{Superato} \\ \hline
		TI44 & Si verifica che InputItemInfoViewImpl interagisca in modo corretto con InputItemInfoViewPresenter & \textcolor{Green}{Implementato} & \textcolor{Green}{Superato} \\ \hline
	\end{longtable}
\end{center}

\subsubsection{Test di Sistema}
Con questa tipologia di test si vuole verificare il corretto funzionamento delle componenti architetturali.\\
I test di sistema saranno descritti nel modo seguente:
\begin{center}
\textbf{TS}[\textit{TipoRequisito}][\textit{ImportanzaRequisito}][\textit{IdRequisito}]
\end{center}
dove:
\begin{itemize}
\item \textbf{TipoRequisito} può assumere valori tra:
\begin{itemize}
\item \textit{F} per i requisiti funzionali;
\item \textit{Q} per i requisiti di qualità;
\item \textit{V} per i requisiti di vincolo;
\item \textit{P} per i requisiti prestazionali.
\end{itemize}
\item \textbf{ImportanzaRequisito} può assumere valori tra:
\begin{itemize}
\item \textit{D} per i requisiti desiderabili;
\item \textit{O} per i requisiti di obbligatori;
\item \textit{F} per i requisiti di facoltativi.
\end{itemize}
\item \textbf{IdRequisito} può assumere un valore gerarchico che identifica il singolo requisito.
\end{itemize}

\subsubsection{Test di Validazione}
I test di validazione hanno lo scopo di accertare che tutte le funzionalità richieste dal proponente siano state soddisfatte.
Per questo motivo, attraverso delle macro azioni, si andrà a simulare il comportamento generale dell'applicativo e dell'utente
che interagisce con esso.
I test di validazione saranno organizzati nel modo seguente:
\begin{center}
\textbf{TV}[\textit{TipoRequisito}][\textit{ImportanzaRequisito}][\textit{IdRequisito}]
\end{center}
dove:
\begin{itemize}
\item
\textbf{TipoRequisito} può assumere valori tra:
	\begin{itemize}
	\item
	\textit{F} per i requisiti funzionali;
	\item
	\textit{Q} per i requisiti di qualità;
	\item
	\textit{V} per i requisiti di vincolo;
	\item
	\textit{P} per i requisiti prestazionali.
	\end{itemize}
\item
\textbf{ImportanzaRequisito} può assumere valori tra:
	\begin{itemize}
	\item
	\textit{D} per i requisiti desiderabili;
	\item
	\textit{O} per i requisiti di obbligatori;
	\item
	\textit{F} per i requisiti di facoltativi.
	\end{itemize}
\item
\textbf{IdRequisito} assume un valore gerarchico che identifica il singolo requisito.
\end{itemize}

\begin{center}
	\begin{longtable}{|c|>{\centering}m{10cm}|c|}\hline
		Id & Descrizione & stato \\ \hline
		TVVO1 & Viene verificato che sia possibile scaricare l'sdk trammite il gestore dei pacchetti npm & Non implementato \\ \hline
		TVFO2 & Viene verificato che l'sdk sia utilizzabile dal applicazione di demo & Non implementato \\ \hline
		TVFO3 & Viene verificato che sia possibile istanziare tutti i widget che sono inclusi nel sdk & Non implementato \\ \hline
		TVFO4 & Viene verificato che sia possibile istanziare tutti i layout che sono inclusi nel sdk & Non implementato \\ \hline
		TVFO5 & Viene verificato che sia possibile istanziare tutte le bolle che sono incluse nel sdk & Non implementato \\ \hline
		TVFO6 & Viene verificato che sia possibile comporre una bolla personalizzata mediante widget e layout & Non implementato \\ \hline
		TVFO7 & Viene verificato che un generico utente possa creare un lista della spesa e pubblicarla al'interno di un canale Rockt.chat & Non implementato \\ \hline
		TVFO8 & Viene verificato che un generico utente possa creare un lista della spesa e pubblicarla ad un singolo utente avente un account sul istanza Rocket.chat su cui è installata l'sdk & Non implementato \\ \hline
		TVFO9 & Viene verificato che se un utente ha i permessi per interagire con una bolla qust'ultimo sia in grado di visualizzare gli elementi e di spuntarli & Non implementato \\ \hline
		TVFO10 & Viene verificato che se un utente inoltra una lista della spesa qust'ultima venga inoltrata in forma testuale bloccando quindi ogni forma di interazione con l'utente & Non implementato \\ \hline
		TVVO11 & Viene  verificato che tutti i file che contengono codice scritto in javascript o coffe script passino i test di linting & Non implementato \\ \hline
	\end{longtable}
\end{center}


\subsection{Strumenti}
Per eseguire i test citati e per ottenere del codice di qualità si è deciso di utilizzare i seguenti strumenti.

\subsubsection{Jenkins}
\termine{Jenkins} è un software che permetterà di automatizzare una serie di operazioni che consentiranno al \termine{team} di monitorare la qualità del software in modo continuo ogni volta che verrà inviata una modifica alla \termine{repository} che conterrà il codice sorgente. \\
Per ottenere questi risultati questo software si appoggerà agli altri strumenti descritti qua di seguito.

\subsubsection{SonarQube}
\termine{SonarQube} è un software che permette di eseguire analisi statica sul codice che verrà inviato alla \termine{repository} fornendo un feedback al programmatore per fargli sapere cosa dovrà migliorare. Inoltre permetterà di controllare che le regole che il \termine{team} si è imposto per la qualità del codice sorgente vengano rispettate.

\subsubsection{Mocha}
\termine{Mocha} è un \termine{framework} che permette di eseguire test di unità e di integrazione andando a creare dei \termine{mock} che simulino il comportamento delle varie parti, con lo scopo di testare ogni dettaglio in modo indipendente dal resto. Ad esso si affiancheranno anche:
\begin{itemize}
\item \textbf{\termine{Chai}}: libreria per formulare asserzioni sull'output atteso;
\item \textbf{\termine{Sinon}}: libreria per simulare delle risposte di richieste a server remoti.
\end{itemize}
\newpage
\section{Pianifazione delle verifiche}

\subsubsection{Verifiche architetturali}
In questa sezione verranno descritte le come e quando effettuare le metriche per valutare l'architettura del software.

\paragraph{Numero di parametri per metodo}
\begin{itemize}
\item \textbf{Chi:} Il \Prog.
\item \textbf{Quando:} Prima della scrittura del metodo.
\item \textbf{Come:} Il \Prog deve controllare prima di progettare il metodo che il numero di parametri soddisfi almeno il livello di accettazione descritto per questa metrica.
\end{itemize}

\paragraph{Livello di stabilità}
\begin{itemize}
\item \textbf{Chi:} il \Ver.
\item \textbf{Quando:} Durante la fase di verifica che precede quella di approvazione.
\item \textbf{Come:} Il \Ver\ deve applicare la formula descritta nelle \NdP.
\end{itemize}
\input{Sezioni/PianificazioneVeriche/VerificheyDocumenti.tex}
\newpage
\subsubsection{Verifiche del software}
In questa sezione verranno descritte come e quando effettuare le misurazioni incentrate sul software.

\paragraph{Complessità ciclomatica}
\begin{itemize}
\item \textbf{Chi:} \termine{SonarQube}.
\item \textbf{Quando:} Dopo ogni push sulla \termine{repository}.
\item \textbf{Come:} Questo test viene effettuato automaticamente da \termine{SonarQube}, se il test non viene superato il codice non verrà integrato nel ramo di produzione.
\end{itemize}

\paragraph{Numero di livelli di annidamento per metodo}
\begin{itemize}
\item \textbf{Chi:} \termine{SonarQube}.
\item \textbf{Quando:} Dopo ogni push sulla \termine{repository}.
\item \textbf{Come:} Questo test viene effettuato automaticamente da \termine{SonarQube}, se il test non viene superato il codice non verrà integrato nel ramo di produzione.
\end{itemize}

\paragraph{Numero di attributi per classe}
\begin{itemize}
\item \textbf{Chi:} \termine{SonarQube}.
\item \textbf{Quando:} Dopo ogni push sulla \termine{repository}.
\item \textbf{Come:} Questo test viene effettuato automaticamente da \termine{SonarQube}, se il test non viene superato il codice non verrà integrato nel ramo di produzione.
\end{itemize}

\paragraph{Copertura dei test}
\begin{itemize}
\item \textbf{Chi:} \termine{SonarQube}.
\item \textbf{Quando:} Dopo ogni push sulla \termine{repository}.
\item \textbf{Come:} Questo test viene effettuato automaticamente da \termine{SonarQube}, se il test non viene superato il codice non verrà integrato nel ramo di produzione.
\end{itemize}

\paragraph{Linee di commento per linee di codice}
\begin{itemize}
\item \textbf{Chi:} Il \Progr.
\item \textbf{Quando:} Dopo la scrittura di ogni metodo o funzione.
\item \textbf{Come:} Il \Progr\ deve verificare che il test soddisfi almeno il livello di accettazione descritto per poter considerare la funzione o il metodo accettabile. Nel caso in cui non sia raggiunto il minimo livello, il \Prog\ dovrà provvedere a trovare frasi più semplici e coese rispetto a quelle in precedenza. 
\end{itemize}

\paragraph{Numero di parametri per metodo}
\begin{itemize}
\item \textbf{Chi:} Il \Prog.
\item \textbf{Quando:} Prima della scrittura del metodo.
\item \textbf{Come:} Il \Prog deve controllare prima di progettare il metodo che il numero di parametri soddisfi almeno il livello di accettazione descritto per questa metrica.
\end{itemize}

\paragraph{Livello di stabilità}
\begin{itemize}
\item \textbf{Chi:} il \Ver.
\item \textbf{Quando:} Durante la fase di verifica che precede quella di approvazione.
\item \textbf{Come:} Il \Ver\ deve applicare la formula descritta nelle \NdP.
\end{itemize}


\end{document}